\scanpage{194}

\uline{Questions}

$S$ schéma de type fini sur $\mathbb Z$, dominant $\mathbb Z$, $S$ intègre pt gén. $\eta$

$X$ schéma lisse de type fini sur $S$.

$E$ module localmt libre sur $X$, avec connexion à courbure nulle relativement à $S$.

Le \inserted{module à} connexion \inserted{$E_{\eta}$ sur $X_{\eta}$} est dit \uline{de niveau zéro} \inserted{ou \uline{intégrable}} \deleted{si l'on peut trouver, sur $X_{\eta}$, un} l'on peut trouver un revêtement fini étale surjectif de $X_{\eta}$ sur lequel $E_{\eta}$ devient localmt trivial, i.e. ``intégrable'' au sens (Zar). \deleted{Une condition}\inserted{\unsure{Il revient|Je veux}} \deleted{nécessaire pour ceci est que} au \unsure{$\simeq$|$=$} de dire que sur une extension alg. close \deleted{\unsure{finie}} $K$ de $k(\eta)$, $E_K$ sur $X_K$ soit associé à des repr. linéaires des groupes fondamentaux des composantes connexes de $E_K$. Et si on choisit $K=\mathbb C$, que les représentations des $\pi_1$ \inserted{transcendants} des composantes \deleted{\unsure{gén.}} connexes, définies par $E_{\mathbb C}$, \deleted{\unsure{soit}} se fassent à travers des groupes finis.

% slanted

\margin{définition a un sens sur n'importe quelle \deleted{\unsure{schéma}} base --- Il faut mieux de dire : localmt \deleted{\unsure{simplemt}} trivial (pour la top. étale) ...}

Cette notion peut se varier aussi pour n'importe quel schéma en groupe de \unsure{t.f.|déf.|dif.|type f.} $G$ sur $S$, on prend sur $X$ un torseur $P$ sous $G_X$, avec connexion à courbure nulle rel/$S$.

La question qui suit pourra se poser aussi pour des tels groupes.

Remarquons que si la connexion de $E$ est intégrable, alors

% archive sheet number: 193
